% !TEX root = TCC - CD e IA.tex
% ---
% RESUMO
% ---
% resumo na língua vernácula (obrigatório)
\setlength{\absparsep}{18pt} % ajusta o espaçamento dos parágrafos do resumo
\begin{resumo}
Este trabalho propõe o desenvolvimento e a validação de uma metodologia híbrida de mineração de texto que serve de motor a uma ferramenta de apoio à decisão para assessorias políticas. A abordagem integra técnicas de Processamento de Linguagem Natural (PLN) --- análise de sentimentos, para identificar polaridades; modelagem de tópicos, para extrair os temas centrais das discussões --- e aprendizado não supervisionado (clusterização), para agrupar publicações e perfis por características semelhantes. Diferentemente de abordagens que se encerram na descrição da opinião pública, os resultados dessas técnicas são organizados em um funil de engajamento inspirado no \textit{marketing} de crescimento (\textit{growth}) e nos construtos de \textit{Consumers' Online Brand-Related Activities} (COBRAs), convertendo a análise em recomendações práticas: o que produzir, o que amplificar e o que abandonar. A metodologia distingue o discurso produzido pelas assessorias --- em legendas e transcrições --- da reação expressa pelo público nos comentários, e é validada no estudo de caso dos vinte e sete governadores brasileiros. Espera-se que a pesquisa contribua com um instrumento analítico replicável para o marketing político digital, superando as limitações de métricas superficiais de engajamento ao qualificá-las pelo sentimento e organizá-las de forma acionável.

\vspace{0.5cm}
\noindent
\textbf{Palavras-chaves}: Marketing Político Digital, Análise de Sentimentos, Modelagem de Tópicos, Engajamento (COBRAs), \textit{Growth}, Ciência de Dados.
\end{resumo}